\chaves{minera\c{c}\~ao de dados educacionais; previs\~ao de desempenho escolar; aprendizado de m\'aquina; risco pedag\'ogico; learning analytics; modelagem temporal tabular.} % Palavras chave separadas por ponto e vírgula

% Não deixe linha em branco após \begin{resumo}
\begin{resumo}
Este Trabalho de Conclus\~ao de Curso apresenta o desenvolvimento de um sistema preditivo voltado ao contexto escolar, com foco na antecipa\c{c}\~ao de dificuldades acad\^emicas e no apoio \`a tomada de decis\~ao de professores, coordena\c{c}\~ao pedag\'ogica e secretaria. O problema central investigado consiste em verificar se dados escolares hist\'oricos, como notas, faltas, pagamentos e informa\c{c}\~oes acad\^emicas, podem ser integrados em uma pipeline anal\'itica capaz de prever a pr\'oxima nota do aluno e identificar sinais de risco pedag\'ogico antes da ocorr\^encia de reprova\c{c}\~ao ou queda severa de desempenho. O projeto teve in\'icio com uma abordagem baseada em grafos, mas evoluiu para uma estrat\'egia tabular temporal ap\'os sucessivas an\'alises indicarem maior ader\^encia entre o problema e modelos supervisionados aplicados a dados longitudinais. A solu\c{c}\~ao final foi estruturada como uma arquitetura reproduz\'ivel centrada no pacote \texttt{school\_predictor}, com separa\c{c}\~ao entre camada privada de prepara\c{c}\~ao e extra\c{c}\~ao do banco, interface p\'ublica baseada em arquivos \texttt{CSV}, engenharia de atributos, valida\c{c}\~ao temporal, sele\c{c}\~ao autom\'atica de modelos, an\'alise de erro e gera\c{c}\~ao de relat\'orios operacionais. Foram comparadas abordagens de baseline, Random Forest e HistGradientBoosting, tanto para regress\~ao da pr\'oxima nota quanto para classifica\c{c}\~ao de risco de nota inferior a 6,0. Os resultados mostraram que, para previs\~ao num\'erica da pr\'oxima nota, o melhor cen\'ario foi obtido com hist\'orico m\'inimo de uma observa\c{c}\~ao por disciplina, alcan\c{c}ando MAE de 0,7601 e superando o baseline simples baseado na \'ultima nota. Para alerta de risco, o melhor cen\'ario ocorreu com hist\'orico m\'inimo de duas observa\c{c}\~oes, alcan\c{c}ando F1 de 0,7437, com recall de 0,8618 no teste temporal. Al\'em da modelagem preditiva, foi desenvolvida uma camada de relat\'orios escolares com sa\'idas espec\'ificas para professor, coordena\c{c}\~ao e secretaria, incluindo um ranking executivo de alunos priorit\'arios. Conclui-se que uma arquitetura anal\'itica temporal, orientada ao uso institucional e acompanhada de mecanismos de interpreta\c{c}\~ao, prioriza\c{c}\~ao e prote\c{c}\~ao de dados, \'e adequada ao contexto escolar e pode apoiar interven\c{c}\~oes mais precoces.
\end{resumo}
